\chapter{Meckel's Diverticulum}

\section*{Definition and Overview}
Meckel's diverticulum is a common congenital gastrointestinal anomaly characterised by a true diverticulum (a blind pouch containing all layers of the intestinal wall) arising from the antimesenteric border of the distal ileum. It is a vestigial remnant of the omphalomesenteric (vitelline) duct, which normally obliterates during the fifth to seventh week of gestation. While the majority of individuals with Meckel's diverticulum remain asymptomatic throughout their lives, it can present with clinical symptoms due to the presence of ectopic tissue (most commonly gastric or pancreatic mucosa), which can cause gastrointestinal haemorrhage, inflammation (diverticulitis), or bowel obstruction.

\section*{Epidemiology}
Meckel's diverticulum is classically described by the ``rule of twos'': it occurs in approximately 2\% of the general population, is located within 2 feet (approximately 60 cm) of the ileocaecal valve, is about 2 inches (5 cm) in length, contains 2 types of ectopic tissue (gastric and pancreatic), presents before the age of 2 years in symptomatic cases, and is twice as common in males as in females for symptomatic presentations. In Australia, it is a significant differential diagnosis in paediatric patients presenting with painless lower gastrointestinal bleeding or acute abdominal pain mimicking appendicitis.

\section*{Aetiology and Pathophysiology}
The development and pathophysiology of Meckel's diverticulum stem from incomplete embryonic involution and the secretion of corrosive fluids.
\begin{itemize}
    \item \textbf{Incomplete vitelline duct obliteration:} Failure of the embryonic connection between the primitive midgut and the yolk sac to atrophy, resulting in a persistent distal ileal pouch.
    \item \textbf{Ectopic tissue migration:} Pluripotent cells within the vitelline duct differentiate into heterotopic mucosa, with gastric mucosa present in over 60\% of cases and pancreatic tissue in about 5\% to 16\%.
    \item \textbf{Acid-induced ulceration:} Ectopic gastric mucosa secretes hydrochloric acid, which corrodes the adjacent, unprotected ileal mucosa, leading to painless ulceration and haemorrhage.
    \item \textbf{Obstruction mechanisms:} The diverticulum can act as a lead point for intussusception (telescoping of the bowel), cause volvulus around a persistent vitelline cord, or lead to internal herniation.
    \item \textbf{Inflammatory diverticulitis:} Obstruction of the diverticulum lumen by a faecolith or foreign body leads to stasis, bacterial overgrowth, and inflammation, mirroring the pathogenesis of acute appendicitis.
\end{itemize}

\section*{Clinical Features}
The clinical presentation of Meckel's diverticulum varies by age, with haemorrhage being more common in children and obstruction or diverticulitis being more common in adults.
\begin{itemize}
    \item \textbf{Painless haematochezia:} Passage of painless, dark red or ``currant-jelly'' stool, often occurring in children under the age of five and sometimes presenting as massive haemorrhage.
    \item \textbf{Abdominal pain:} Colicky or constant abdominal pain, localising to the periumbilical or right lower quadrant, often accompanied by tenderness and guarding.
    \item \textbf{Bowel obstruction:} Nausea, bilious vomiting, abdominal distension, and inability to pass gas or stool, secondary to intussusception or volvulus.
    \item \textbf{Signs of anaemia:} Pallor, fatigue, tachycardia, and lethargy in children with chronic, low-grade bleeding.
    \item \textbf{Peritonitis:} Severe abdominal pain, rigidity, and high fever if the diverticulum perforates due to inflammation or ulceration.
\end{itemize}

\section*{Differential Diagnosis}
Meckel's diverticulum must be differentiated from other common causes of acute abdomen and gastrointestinal bleeding.
\begin{itemize}
    \item \textbf{Acute appendicitis:} Presents with fever, right lower quadrant pain, and leukocytosis, and is clinically indistinguishable from acute Meckel's diverticulitis without surgical exploration or imaging.
    \item \textbf{Inflammatory bowel disease (IBD):} Conditions like Crohn's disease can present with terminal ileitis, abdominal pain, and gastrointestinal bleeding, but usually show a more chronic, systemic course.
    \item \textbf{Intussusception (idiopathic):} Telescoping of the bowel in infants, often without a specific lead point, presenting with severe colicky pain and currant-jelly stools.
    \item \textbf{Juvenile polyps:} Benign mucosal outgrowths in the colon or rectum that cause painless rectal bleeding in children, differentiated by colonoscopy.
    \item \textbf{Gastroenteritis:} Viral or bacterial bowel infection presenting with diarrhoea, vomiting, and abdominal cramps, but lacking localised peritonitis or massive painless bleeding.
\end{itemize}

\section*{When to Seek Medical Care}
Parents or individuals should seek immediate medical attention if they notice blood in the stool or experience severe, progressive abdominal pain.

\textbf{Call 000 (or local emergency services) for:}
\begin{itemize}
    \item Large volume rectal bleeding or signs of hypovolaemic shock (such as extreme pallor, rapid breathing, cold and clammy skin, or loss of consciousness).
    \item Severe, sudden abdominal pain with a rigid, board-like abdomen and high fever, indicating bowel perforation or peritonitis.
    \item Uncontrolled, persistent vomiting accompanied by severe abdominal swelling.
\end{itemize}

\section*{Diagnosis and Investigations}
Diagnosing a Meckel's diverticulum can be challenging, as it is often not visible on routine imaging unless inflamed.
\begin{itemize}
    \item \textbf{Meckel's scan (Technetium-99m pertechnetate scintigraphy):} The diagnostic test of choice in children, where the radiotracer is taken up by ectopic gastric mucosa, demonstrating a ``hot spot'' in the right lower quadrant.
    \item \textbf{Abdominal ultrasound:} High-resolution ultrasound can identify an inflamed diverticulum, an intussusception, or a fluid-filled pouch in the lower abdomen.
    \item \textbf{CT scan of the abdomen:} Useful in adults and acute settings to detect diverticulitis, enteroliths, or signs of bowel obstruction and perforation.
    \item \textbf{Angiography or RBC scan:} In cases of active, rapid bleeding, to identify the site of contrast extravasation or pooling of labelled red blood cells.
    \item \textbf{Diagnostic laparoscopy:} Often the definitive diagnostic and therapeutic tool, especially when patients present with an acute abdomen and appendicitis is suspected but a normal appendix is found.
\end{itemize}

\section*{Management}
Management is surgical excision for symptomatic diverticula, while the management of asymptomatic diverticula discovered incidentally remains controversial.
\begin{itemize}
    \item \textbf{Surgical resection:} Definitive treatment involves laparoscopic or open diverticulectomy (removal of the diverticulum) or segmental ileal resection (preferred if there is ulceration or bleeding in the adjacent ileum).
    \item \textbf{Haemodynamic stabilisation:} Intravenous fluid resuscitation and blood transfusions as required for patients presenting with haemorrhage or shock.
    \item \textbf{Antibiotic therapy:} Intravenous broad-spectrum antibiotics (e.g.\ ampicillin, gentamicin, and metronidazole) for patients with diverticulitis or perforation.
    \item \textbf{Nasogastric decompression:} Placement of a nasogastric tube to decompress the stomach and bowel in patients presenting with bowel obstruction.
    \item \textbf{Incidental management:} Asymptomatic diverticula found during abdominal surgery are often resected in patients under 40, males, or if the diverticulum is >2 cm or has palpable ectopic tissue, to prevent future complications.
\end{itemize}

\section*{Complications}
Untreated Meckel's diverticulum can lead to severe, life-threatening complications.
\begin{itemize}
    \item \textbf{Bowel perforation:} Erosion of the diverticulum or adjacent ileal wall by acid or inflammation, leading to peritonitis and sepsis.
    \item \textbf{Severe haemorrhage:} Massive lower gastrointestinal bleeding leading to acute anaemia, hypovolaemic shock, and cardiovascular instability.
    \item \textbf{Gangrene and necrosis:} Strangulation of the diverticulum or loop of bowel due to volvulus or intussusception, requiring bowel resection.
    \item \textbf{Chronic iron-deficiency anaemia:} Low-grade, occult bleeding from chronic ulceration causing long-term iron depletion.
\end{itemize}

\section*{Prognosis and Outcomes}
The prognosis for symptomatic Meckel's diverticulum following surgical resection is excellent. Most patients achieve complete recovery and do not experience long-term digestive or nutritional complications. The risk of surgical complications (such as wound infection, bleeding, or postoperative adhesions) is low, particularly when performed laparoscopically. Following resection, there is zero risk of recurrence.

\section*{Prevention and Self-Care}
Because Meckel's diverticulum is a congenital developmental anomaly, there are no known preventive measures or home remedies.
\begin{itemize}
    \item \textbf{Monitor paediatric symptoms:} Parents should remain vigilant for unexplained blood in their child's stool or signs of recurrent abdominal pain.
    \item \textbf{Follow postoperative care:} Post-surgical patients should adhere to wound care instructions, take prescribed analgesics, and gradually reintroduce solid foods.
    \item \textbf{Recognise bowel obstruction signs:} Be aware of signs of postoperative bowel obstruction, such as abdominal pain, vomiting, or failure to pass flatus, and seek prompt care if they occur.
\end{itemize}

\section*{Sources and Further Reading}
The following reputable organisations provide further information on Meckel's diverticulum:
\begin{itemize}
    \item Sydney Children's Hospitals Network. \textit{Fact sheets and clinical guidelines on paediatric surgical conditions.} https://www.schn.health.nsw.gov.au/
    \item Healthdirect Australia. \textit{Overview of Meckel's diverticulum, symptoms, and surgical management.} https://www.healthdirect.gov.au/meckels-diverticulum
    \item Mayo Clinic. \textit{Clinical guide to congenital gastrointestinal anomalies and Meckel's diverticulum.} https://www.mayoclinic.org/
    \item NHS Great Ormond Street Hospital. \textit{Patient information and treatment pathways for Meckel's diverticulum in children.} https://www.gosh.nhs.uk/
    \item MSD Manual (Professional Version). \textit{Detailed pathophysiology, diagnosis, and management of Meckel's diverticulum.} https://www.msdmanuals.com/
\end{itemize}
